\chapter*{Synthèse}
\addcontentsline{toc}{chapter}{Synthèse}

L'évaluation des engagements financiers sous le référentiel Solvabilité II exige des projections stochastiques lourdes, générant des temps de calcul souvent incompatibles avec les contraintes opérationnelles des organismes d'assurance. Cette étude s'attache à résoudre cette problématique de dimensionnalité en explorant diverses techniques d'agrégation appliquées à un portefeuille de passifs d'assurance vie. L'objectif consiste à agréger l'information des contrats tout en préservant la fidélité des indicateurs de risque réglementaires, notamment le \textit{Best Estimate} (BE) et le capital de solvabilité requis (SCR).

Face à la volumétrie massive des portefeuilles d'assurance vie, la méthode usuelle d'agrégation repose sur la constitution de \textit{Model Points} (MP). Ces portefeuilles agrégés regroupent des polices présentant des profils de risque similaires. L'enjeu technique majeur réside dans le calibrage de cette compression : une agrégation excessive engendre un lissage des garanties, altérant le calcul du BE et du SCR. Pour s'aligner sur les standards opérationnels du marché, l'étude impose une contrainte de compression stricte, limitant la taille du portefeuille à une volumétrie comprise entre 800 et 1,200 \textit{Model Points}.

Pour mener cette analyse, un portefeuille de référence synthétique de 1,000,000 de polices a été généré, reflétant fidèlement la structure du marché français. Ce générateur stochastique s'appuie sur une approche par générations, assurant la cohérence temporelle et financière des caractéristiques techniques simulées. Les distributions statistiques retenues ont été calibrées sur des données publiques empiriques :
\begin{itemize}
\item \textbf{Âge de l'assuré :} Loi Bêta, reproduisant l'asymétrie démographique des détenteurs d'assurance vie ;
\item \textbf{Sexe :} Loi de Bernoulli, paramétrée selon la répartition sectorielle observée ;
\item \textbf{Ancienneté du contrat :} Modèle de mélange gaussien (\textit{Gaussian Mixture Model} - GMM), sélectionné pour capter les différentes vagues historiques de collecte ;
\item \textbf{Provisions mathématiques (PM) :} Loi Log-Normale, capturant l'asymétrie de la distribution des richesses et la concentration des encours sur les contrats à forts capitaux ;
\item \textbf{Caractéristiques techniques conditionnelles :} Le taux minimum garanti (TMG) et la part d'unités de compte (UC) sont conditionnés à l'année de souscription pour refléter l'évolution de l'environnement économique.
\end{itemize}

Afin de confronter les différentes logiques de réduction de dimension, une agrégation stricte à iso-caractéristiques sert de ligne de base (\textit{Baseline}). N'induisant aucune perte d'information, elle est comparée à cinq approches distinctes :
\begin{itemize}
    \item \textbf{Le découpage en classes (\textit{Banding}) :} Méthode déterministe usuelle segmentant le portefeuille selon des maillages prédéfinis appliqués aux variables clés (âge, ancienneté, TMG) ;
    \item \textbf{L'algorithme \textit{K-Means} :} Méthode d'apprentissage non supervisé procédant par partitionnement spatial afin de minimiser la variance intra-classe des attributs statiques ;
    \item \textbf{Les algorithmes \textit{DBSCAN} et \textit{HDBSCAN} :} Approches spatiales basées sur la densité locale, permettant d'isoler les profils atypiques via la gestion du bruit. Le traitement des variables discrètes requiert l'implémentation préalable d'un bruit blanc uniforme (\textit{Dithering}) ;
    \item \textbf{Le \textit{Cash-Flow Matching} :} Méthode d'agrégation dynamique s'appliquant non pas aux caractéristiques statiques à la date d'arrêté, mais à la morphologie de l'écoulement futur des flux de trésorerie. L'algorithme \textit{K-Means} est employé pour regrouper les polices présentant des trajectoires de flux similaires ;
    \item \textbf{La calibration \textit{a posteriori} :} Approche consistant à regrouper les polices selon leurs caractéristiques financières. Des lois de sortie sont ensuite calibrées pour chaque regroupement afin de restituer le BE du portefeuille de référence à l'identique sur une projection en passif seul déterministe.
\end{itemize}

Le protocole de sélection des hyperparamètres repose sur une recherche exhaustive (\textit{Grid Search}), visant à identifier l'optimum de Pareto pour chaque algorithme. Il s'agit d'atteindre le meilleur compromis entre la contrainte opérationnelle de compression imposée et la minimisation de l'erreur sur le \textit{Best Estimate} (BE), évalué dans un premier temps à l'aide d'un modèle en passif seul déterministe. 

Le tableau \ref{tab:synthese_calibration_finale_synthese} présente les résultats obtenus sur le modèle passif seul. Il montre que toutes les méthodes de regroupement s'avèrent très précises dans ce cadre de base, avec des erreurs négligeables par rapport au portefeuille initial complet.

\begin{table}[H]
    \centering
    \renewcommand{\arraystretch}{1.4}
    \begin{tabular}{|l|r|r|r|}
        \hline
        \textbf{Méthode d'agrégation} & \textbf{Nb MP} & \textbf{BE Passif Seul} & \textbf{Erreur Passif Seul} \\
        \hline
        \textit{Baseline} (Référence) & 26\,272 & 263{,}95 Md\euro{} & \textit{Réf.} \\
        \textit{Banding} Dynamique & 1\,172 & 264{,}03 Md\euro{} & $+0{,}0292\,\%$ \\
        \textit{K-Means} Robuste & 1\,039 & 263{,}96 Md\euro{} & $+0{,}0035\,\%$ \\
        \textit{DBSCAN} & 1\,047 & 263{,}98 Md\euro{} & $+0{,}0124\,\%$ \\
        \textit{HDBSCAN} & 967 & 264{,}06 Md\euro{} & $+0{,}0398\,\%$ \\
        \textit{Cash-Flow Matching} & 1\,000 & 263{,}94 Md\euro{} & $-0{,}0049\,\%$ \\
        Calibration \textit{a posteriori} & 222 & 263{,}95 Md\euro{} & $0{,}0000\,\%$ \\
        \hline
    \end{tabular}
    \caption{Évaluation déterministe en passif seul et erreurs d'agrégation associées.}
    \label{tab:synthese_calibration_finale_synthese}
\end{table}

La forte compression de la volumétrie, divisée par plus de vingt-cinq par rapport au portefeuille de référence, s'opère sans déformation significative du BE, l'erreur maximale plafonnant à $+0{,}0398\,\%$. Par ailleurs, la restitution parfaite de la calibration \textit{a posteriori} découle mécaniquement de sa construction, les lois de sortie étant spécifiquement ajustées pour répliquer l'engagement déterministe de référence.

Avant d'éprouver la robustesse de ces modèles face aux chocs réglementaires, l'analyse des temps d'exécution (Tableau \ref{tab:runtimes_synthese}) chiffre le gain de performance obtenu lors de la transposition en environnement ALM stochastique complet.

\begin{table}[H]
    \centering
    \renewcommand{\arraystretch}{1.4}
    \begin{tabular}{|l|r|r|r|}
        \hline
        \textbf{Méthode d'agrégation} & \textbf{Volumétrie (MP)} & \textbf{Temps d'exécution} & \textbf{Facteur de réduction} \\
        \hline
        \textit{Baseline} (Référence) & 26\,272 & $\approx$ 8 h 00 (28\,783 s) & \textit{Réf.} \\
        \textit{HDBSCAN} & 967 & $\approx$ 10 min 00 (600 s) & $\div 48$ \\
        \textit{Cash-Flow Matching} & 1\,000 & $\approx$ 10 min 27 (627 s) & $\div 46$ \\
        \textit{K-Means} Robuste & 1\,039 & $\approx$ 11 min 11 (671 s) & $\div 43$ \\
        \textit{Banding} Dynamique & 1\,172 & $\approx$ 11 min 51 (711 s) & $\div 40$ \\
        Calibration \textit{a posteriori} & 222 & $\approx$ 2 min 38 (158 s) & $\div 182$ \\
        \hline
    \end{tabular}
    \caption{Temps de calcul comparés en projection ALM stochastique (5\,000 scénarios).}
    \label{tab:runtimes_synthese}
\end{table}

La réduction de la volumétrie divise les temps de traitement par un facteur 46 à 1000 MP, ramenant une projection exhaustive de huit heures à une dizaine de minutes. Ce gain de performance permet la multiplication des calculs de sensibilités sur des infrastructures standards.

La validation de l'agrégation requiert l'analyse du bilan sous le scénario stochastique central. Le modèle ALM intègre les interactions non linéaires entre les rendements financiers et le comportement des assurés via la mécanique de la participation aux bénéfices. Le tableau \ref{tab:resultats_central_alm} expose l'évaluation du BE et du SCR dans ce cadre d'étude.

\begin{table}[H]
    \centering
    \renewcommand{\arraystretch}{1.4}
    \begin{tabular}{|l|r|r|r|r|}
        \hline
        \textbf{Méthode d'agrégation} & \textbf{BE Central} & \textbf{Erreur BE} & \textbf{SCR Global} & \textbf{Erreur SCR} \\
        \hline
        \textit{Baseline} (Référence) & 209{,}93 Md\euro{} & \textit{Réf.} & 26{,}91 Md\euro{} & \textit{Réf.} \\
        \textit{Banding} Dynamique & 209{,}02 Md\euro{} & $-0{,}43\,\%$ & 27{,}47 Md\euro{} & $+2{,}10\,\%$ \\
        \textit{HDBSCAN} & 209{,}05 Md\euro{} & $-0{,}42\,\%$ & 27{,}47 Md\euro{} & $+2{,}10\,\%$ \\
        \textit{K-Means} Robuste & 208{,}95 Md\euro{} & $-0{,}46\,\%$ & 27{,}49 Md\euro{} & $+2{,}17\,\%$ \\
        \textit{Cash-Flow Matching} & 208{,}41 Md\euro{} & $-0{,}72\,\%$ & 27{,}70 Md\euro{} & $+2{,}94\,\%$ \\
        Calibration \textit{a posteriori} & 198{,}31 Md\euro{} & $-5{,}53\,\%$ & 19{,}27 Md\euro{} & $-28{,}38\,\%$ \\
        \hline
    \end{tabular}
    \caption{Évaluation globale des engagements et du SCR en projection ALM stochastique.}
    \label{tab:resultats_central_alm}
\end{table}

Les algorithmes spatiaux (\textit{HDBSCAN}, \textit{K-Means}) et le découpage par classes (\textit{Banding}) démontrent une grande précision, limitant la déviation du SCR global sous le seuil de $2{,}20\,\%$. Afin d'isoler l'origine des biais d'agrégation, la décomposition de l'erreur sur les principaux modules de risque est présentée dans le tableau \ref{tab:resultats_modules}.

\begin{table}[H]
    \centering
    \renewcommand{\arraystretch}{1.4}
    \begin{tabular}{|l|r|r|r|}
        \hline
        \textbf{Méthode d'agrégation} & \textbf{Erreur SCR Marché} & \textbf{Erreur SCR Vie} & \textbf{Erreur SCR Global} \\
        \hline
        \textit{Banding} Dynamique & $+1{,}89\,\%$ & $+3{,}55\,\%$ & $+2{,}10\,\%$ \\
        \textit{HDBSCAN} & $+1{,}87\,\%$ & $+3{,}71\,\%$ & $+2{,}10\,\%$ \\
        \textit{K-Means} Robuste & $+1{,}92\,\%$ & $+3{,}87\,\%$ & $+2{,}17\,\%$ \\
        \textit{Cash-Flow Matching} & $+2{,}42\,\%$ & $+6{,}50\,\%$ & $+2{,}94\,\%$ \\
        Calibration \textit{a posteriori} & $-52{,}82\,\%$ & $+83{,}88\,\%$ & $-28{,}38\,\%$ \\
        \hline
    \end{tabular}
    \caption{Décomposition de l'erreur sur les principaux modules de risque.}
    \label{tab:resultats_modules}
\end{table}

L'approche dynamique du \textit{Cash-Flow Matching} affiche une surestimation marquée sur le SCR Vie ($+6{,}50\,\%$). Fondée exclusivement sur la morphologie centrale des flux, cette méthode occulte l'hétérogénéité de l'ancienneté fiscale au sein des regroupements, générant une sur-réaction mécanique aux hypothèses de rachats structurels. Enfin, la méthode de calibration \textit{a posteriori} s'avère inopérante en environnement stochastique dû au fait que des lois comportementales figées n'arrivent pas à répliquer la réalité du portefeuille dans un environnement stochastique.

L'évaluation de la résilience des modèles nécessite l'application de sensibilités sur l'environnement économique. Le portefeuille est ainsi soumis à deux déformations parallèles de la courbe des taux sans risque publiée par l'EIOPA, s'élevant respectivement à $-100$ et $+100$ points de base. Cette étape éprouve la résistance des méthodes de regroupement face à des environnements économiques différents (Tableau \ref{tab:sensibilites_scr_synthese}).

\begin{table}[H]
    \centering
    \renewcommand{\arraystretch}{1.4}
    \resizebox{\textwidth}{!}{
    \begin{tabular}{|l|r|r|r|r|}
        \hline
        \multirow{2}{*}{\textbf{Méthode d'agrégation}} & \multicolumn{2}{c|}{\textbf{Baisse (-100 bps)}} & \multicolumn{2}{c|}{\textbf{Hausse (+100 bps)}} \\
        \cline{2-5}
        & \textbf{Erreur SCR Rachats} & \textbf{Erreur SCR Global} & \textbf{Erreur SCR Rachats} & \textbf{Erreur SCR Global} \\
        \hline
        \textit{Banding} Dynamique & $+7{,}55\,\%$ & $+1{,}17\,\%$ & $+4{,}55\,\%$ & $+1{,}23\,\%$ \\
        \textit{HDBSCAN} & $+8{,}22\,\%$ & $+1{,}16\,\%$ & $+4{,}68\,\%$ & $+1{,}24\,\%$ \\
        \textit{K-Means} Robuste & $+11{,}56\,\%$ & $+1{,}21\,\%$ & $+4{,}72\,\%$ & $+1{,}25\,\%$ \\
        \textit{Cash-Flow Matching} & $+1{,}11\,\%$ & $+1{,}39\,\%$ & $+6{,}23\,\%$ & $+1{,}27\,\%$ \\
        Calibration \textit{a posteriori} & $+1779{,}84\,\%$ & $-19{,}19\,\%$ & $+0{,}47\,\%$ & $-37{,}16\,\%$ \\
        \hline
    \end{tabular}
    }
    \caption{Erreurs relatives sur le SCR Global et le SCR Rachats selon la déformation de la courbe des taux.}
    \label{tab:sensibilites_scr_synthese}
\end{table}

L'exigence en capital au titre du risque de rachat présente une forte asymétrie en fonction de l'environnement financier. Dans un scénario de taux bas, le rendement des actifs est absorbé par le coût des taux minimums garantis (TMG), la rentabilité future du portefeuille est donc restreinte. L'application d'un choc réglementaire de rachat (choc à la hausse ou rachat de masse) ne diminue pas beaucoup la marge et donc n'affecte pas significativement le SCR. En revanche, une hausse des taux d'intérêt restaure un différentiel de rendement très positif. Le choc de rachat vient alors empêcher une marge importante dans le futur. C'est la matérialisation de ce coût d'opportunité financier qui justifie la forte augmentation du SCR Rachats à la hausse.

Les algorithmes spatiaux, en préservant la granularité des taux garantis et des variables statiques pilotant les lois de rachats structurels (telle que l'ancienneté), restituent cette perte de valeur avec une erreur contenue autour de $+4{,}70\,\%$ à la hausse. À l'inverse, la méthode du \textit{Cash-Flow Matching} affiche une déviation supérieure ($+6{,}23\,\%$). En lissant les caractéristiques individuelles au profit d'une agrégation fondée uniquement sur la morphologie des flux centraux, cette approche altère l'application des taux de rachats structurels de base, générant une sur-réaction du portefeuille face aux chocs réglementaires.

En conclusion, cette étude valide l'usage des algorithmes d'apprentissage non supervisé pour l'agrégation dans le cadre de la modélisation ALM, tout en apportant un éclairage critique sur les pratiques de place :

\begin{itemize}
    \item \textbf{La robustesse de l'approche empirique :} La méthode par découpage en tranches (\textit{Banding}) démontre une résilience marquée face aux chocs asymétriques. Elle confirme que les variables statiques, telles que le TMG et l'ancienneté fiscale, demeurent des axes de regroupement structurants. Toutefois, la rigidité de son maillage contraint mécaniquement ses capacités de compression maximale.
    \item \textbf{L'optimum industriel par le \textit{Machine Learning} :} Pour repousser les limites de la compression, les algorithmes non supervisés se démarquent. Si \textit{HDBSCAN} affiche une précision théorique remarquable grâce à sa gestion fine de la densité locale, la nature de ses hyperparamètres rend le ciblage d'une volumétrie précise particulièrement complexe. L'algorithme \textit{K-Means} s'impose ainsi comme le meilleur compromis opérationnel : il assure un contrôle direct sur le nombre de \textit{Model Points} générés tout en restituant la déformation de l'exigence en capital avec une grande fidélité.
    \item \textbf{Les limites des approches dynamiques :} La méthode du \textit{Cash-Flow Matching} s'avère efficiente sur les projections centrales, mais son lissage des caractéristiques individuelles la pénalise lors de l'évaluation des chocs réglementaires. Cette déviation affecte non seulement le risque de rachat, mais s'étend mécaniquement à l'ensemble des modules de risque sensibles aux comportements des assurés.
\end{itemize}

La portée de ces conclusions reste toutefois soumise aux hypothèses de l'étude. D'une part, l'utilisation d'un générateur stochastique, bien que fondé sur des données empiriques, introduit nécessairement une simplification par rapport à la complexité d'un portefeuille réel. D'autre part, les tests de sensibilités se sont limités au seul risque de taux. Une évaluation intégrant d'autres changements dans les hypothèses de projection permettrait d'éprouver plus largement les méthodes d'agrégation. De plus, l'enrichissement du modèle ALM, via l'implémentation d'un algorithme de participation aux bénéfices plus sophistiqué et l'introduction d'une loi de rachats dynamiques, constituerait un axe de recherche particulièrement pertinent pour tester la robustesse des regroupements face à des comportements fortement non linéaires.

Ces travaux ouvrent plusieurs axes d'amélioration pour la modélisation ALM. La mécanique de l'algorithme \textit{HDBSCAN} offre une voie de recherche particulièrement pertinente. Plutôt que d'imposer un nombre défini de \textit{Model Points}, cet algorithme permet de déduire la volumétrie optimale directement à partir de la densité des données. L'apprentissage non supervisé permet ainsi de dépasser les cibles de compression, au profit d'une modélisation pilotée par la nature même des engagements financiers. D'une autre part, les approches dynamiques pourraient être enrichies et calibrées sur des projections ALM afin de mieux capturer les interactions Actif-Passif et donc de limiter les biais d'agrégation lors de l'évaluation des chocs réglementaires. Enfin, la validation de ces méthodes sur des portefeuilles réels permettrait de confirmer la pertinence opérationnelle de ces techniques d'agrégation dans un contexte industriel.